\section{New runtime and compiler-interface results}
\label{sec:runtime-new}

This investigation found additional freedom in the interface between the
literal register-machine runtime and its compressed grammar. It did not
construct a verified machine below 245 states. In particular, the negative
search results below do not constitute a lower bound on Turing-machine size.

\subsection{Observed incoming bits are not semantic contracts}

The source extractor associates each closed atomic block with the final
incoming bits observed at its original occurrences. Those masks are useful
conservative restrictions, but an unobserved incoming bit can be valid.
The accompanying increment-orientation lemma gives an ordinary mathematical
proof, using the inherited arbitrary-register mark and increment contracts,
that existing states 39, 72, and 93 admit both incoming bits. These entries
implement increments of $X,H,Y$, respectively. Their original grammar symbols
are $h,m,r$, and their fine-alphabet tokens are $A,T,F$. No literal transition
changes and no extra state is charged. State 70, already used on both sides,
and existing states 169 and 201 satisfy the same lemma.

The exact local row premises are checked by
\texttt{check\_increment\_orientation\_rows.py}. This check is not a proof-kernel
certificate for the induction hypotheses. The mathematical argument is in
\texttt{increment\_orientation\_lemma.tex}.

\subsection{A wider symbolic survey of concrete control contexts}

We then preserved the inherited affine symbolic executor unchanged and
replaced only the external decoder by a transparent virtual decoder. Internal
atomic syntax is reconstructed using the actual Routing245 reassociation.
The experiment quantifies register values, nonempty finite spare suffixes,
and header slack through the inherited induction contracts. Its external
pointer prefixes are concrete, so it must not be promoted into a theorem
about all possible control contexts.

For each unambiguous physical atomic entry except the complete-program entry
150, the exhaustive short-prefix survey considers every nonempty binary
outer prefix of length at most seven with no seven consecutive zeros.
There are 126 ending in zero and 127 ending in one. Primitive paths whose
concatenation violates this scanner precondition are excluded. The executor
checked 42,796 primitive cases before the configured first failure in each
rejected orientation. Fourteen previously one-sided entries passed both
orientations throughout the tested contexts:
\[
17,36,39,67,72,77,93,110,122,125,127,146,149,204.
\]
The three increments above have the separate general argument. For the
remaining entries, these survey results are search guidance: a concrete
candidate must replay every one of its actual terminal paths on its own
literal table. The same test also rejects state 81 in a synthetic context
matching its originally observed bit, demonstrating why an observed mask
alone is not a sufficient closed-block specification.

This boundary was tested constructively during the investigation. A new
128-rule, fully bootstrap-unlocked grammar used three additional orientations
(states 77, 67, and 127). Its source bisimulation passed, but replay on its
compiled literal table failed at pointer $001010$. Its blank execution
also reaches state 42 after 51 steps with head 5 and ones exactly
$\{-3,0,2,3\}$; the row $42/0=1R42$ then writes forever along a blank right
ray. This is a rejected candidate, not an independence machine. The example
does not isolate terminal orientation as the cause: further reassociation
inside the physical atoms and the unlocked bootstrap are additional changed
conditions. It does establish that the finite baseline-atom survey cannot
replace the full candidate proof.

State 150, the complete changed program, was tested separately on three
outer prefixes per orientation. It passed 25,020 primitive obligations per
orientation. This smaller finite set is not included in the fourteen-entry
short-prefix claim.

\subsection{A suffix-sensitive increment entry}

State 18, increment of $N$, is different from the generic increment family.
Its zero row enters state 27 one cell to the left. If the pointer ends in one,
the row $27/1=1R39$ restores the generic pulse. If the pointer ends in $10$,
the rows
\[
27/0=0L32,\quad32/1=1R46,\quad46/0=0R39
\]
do the same. A singleton zero pointer also uses the preceding marker one.
Thus the additional safe left orientation has a two-bit context condition:
the outer pointer must end in $10$, with the marker handling the singleton
case. By contrast, a suffix $00$ routes through
\[
32/0=0R45,\quad45/0=0R21,\quad21/0=0R30,
\]
and the third pulse uses $30/1=0R95$, the decrement backend, instead of the
increment backend 73. This is a concrete reason to retain a context
restriction, even after relaxing the other increments.

A second survey considered all 506 legal outer prefixes of lengths one
through eight for states 18, 51, 68, 144, 145, and 148. Its affine outcomes
match exactly the predicates ``not ending in $00$'' for 18, ``ending in zero''
for 51, 144, and 148, and ``ending in one'' for 68 and 145. This is finite
control evidence for those predicates, not a completeness theorem.
A grammar compiler could propagate context-dependent word masks: for
example, a word having state 18 on its left edge may be admissible as a right
child even when it is inadmissible as a left child. That is a distinct next
optimization from simply declaring every terminal two-sided.

\subsection{Simultaneous deletion with two independently replaced incoming actions}

Earlier independent-splice screening can miss a compensating pair of edits
when each edit fails on its own. We therefore tested the two actions
simultaneously. For every eligible state with exactly two external incoming
rows, delete that state and replace both incoming actions independently by
any write bit, direction, and remaining working-state target. There are
244 remaining targets and four write--direction combinations, so each of
the 54 deleted-state choices defines
\[
(4\cdot244)^2=952,576
\]
action pairs, a combined family of 51,439,104.

The finite fixtures are the inherited 8,385 exact whole-frame representative
cases in \nolinkurl{candidate245-literal-cases.txt.gz}. They use header distance
181 and 167 spare blocks of value one. The unchanged baseline passes every
case and exercises all 490 literal rows; the maximum baseline run length
is 2,686 steps. The search uses a 20,000-step limit and a 2,048-cell window,
with exact final tape and head comparison at dispatcher 96. Halt is treated
separately. It does not assume that the fixtures are blank-reachable.

Sound prefix pruning rejects an entire first-action branch only if the
chosen witness fails before encountering the second edited row. Otherwise
the second actions are enumerated from the suspended configuration. The
completed run executes 9,555,040 joint branches; 2,703 pass its initial
small witness set, and none pass all affected fixtures. Unaffected baseline
fixtures have unchanged executions. These counts and all 54 state choices
are recorded in \texttt{coupled-complete.jsonl}. The family excludes adding
new states, changing the entry invariant, or changing any other row; targets
are remaining working states, not the halt state. The bounded result says
only that this specified family produced no accepted candidate.

\subsection{Joint grammar and scanner cost}

The scanner currently needs eight zero-detection states because a valid
pointer can contain six zeros and an intermediate phase contributes one
additional zero. A previously audited 116-state framework is available
when the source pointer bound is reduced to five. Consequently, two paths
would each reach 244 states:
\[
117+127=244,\qquad116+128=244.
\]
The latter requires a genuine zero-five grammar and its compatible runtime
and initializer, not deletion of a scanner state from the current grammar.
Reducing total derivation depth alone does not save a scanner state: total
depth and maximum consecutive-zero length are different resources.

A new joint objective searches
\[
\#\mathrm{grammar\ rules}
  +\mathbf1\{\mathrm{maximum\ zero\ run}=6\},
\]
subject to zero-run bound six and depth bound 64. Cost at most 128 is the
prospective threshold, before whole-frame and blank-start verification.
This runtime lane uses seed 181804 for 600 seconds with the fourteen-mask,
bootstrap-unlocked source supplied by the grammar investigation. The exact
final metadata and outcome are recorded in \texttt{joint-cost-run.json};
the grammar is an abstract search artifact until compiled and replayed.

\paragraph{Verification boundary.}
Neither the finite deletion search, the control-prefix survey, nor an
abstract grammar count proves a new consistency machine. The positive
reusable result here is the increment interface lemma, together with
explicit additional search freedom that can be discharged on a complete
candidate. No new Lean or Rocq endpoint is claimed.
